\documentclass[10pt,a4paper]{amsart}
\usepackage{amsmath,amssymb,amsthm,graphicx}
\usepackage{verbatim}
\newenvironment{reference}{\comment}{\endcomment}
%\newenvironment{reference}{}{}
\newenvironment{slogan}{\comment}{\endcomment}
\newenvironment{history}{\comment}{\endcomment}

% For commutative diagrams we use Xy-pic
\usepackage[all]{xy}

% We use 2cell for 2-commutative diagrams.
\xyoption{2cell}
% Theorem environments.
%
\theoremstyle{plain}
\newtheorem{theorem}[subsection]{Theorem}
\newtheorem{proposition}[subsection]{Proposition}
\newtheorem{lemma}[subsection]{Lemma}

\theoremstyle{definition}
\newtheorem{definition}[subsection]{Definition}
\newtheorem{example}[subsection]{Example}
\newtheorem{exercise}[subsection]{Exercise}
\newtheorem{situation}[subsection]{Situation}

\theoremstyle{remark}
\newtheorem{remark}[subsection]{Remark}
\newtheorem{remarks}[subsection]{Remarks}

\numberwithin{equation}{subsection}

% Macros
%
\def\lim{\mathop{\mathrm{lim}}\nolimits}
\def\colim{\mathop{\mathrm{colim}}\nolimits}
\def\Spec{\mathop{\mathrm{Spec}}}
\def\Hom{\mathop{\mathrm{Hom}}\nolimits}
\def\Ext{\mathop{\mathrm{Ext}}\nolimits}
\def\SheafHom{\mathop{\mathcal{H}\!\mathit{om}}\nolimits}
\def\SheafExt{\mathop{\mathcal{E}\!\mathit{xt}}\nolimits}
\def\Sch{\mathit{Sch}}
\def\Mor{\mathop{\mathrm{Mor}}\nolimits}
\def\Ob{\mathop{\mathrm{Ob}}\nolimits}
\def\Sh{\mathop{\mathit{Sh}}\nolimits}
\def\NL{\mathop{N\!L}\nolimits}
\def\CH{\mathop{\mathrm{CH}}\nolimits}
\def\proetale{{pro\text{-}\acute{e}tale}}
\def\etale{{\acute{e}tale}}
\def\QCoh{\mathit{QCoh}}
\def\Ker{\mathop{\mathrm{Ker}}}
\def\Im{\mathop{\mathrm{Im}}}
\def\Coker{\mathop{\mathrm{Coker}}}
\def\Coim{\mathop{\mathrm{Coim}}}

% Boxtimes
%
\DeclareMathSymbol{\boxtimes}{\mathbin}{AMSa}{"02}

%
% Macros for moduli stacks/spaces
%
\def\QCohstack{\mathcal{QC}\!\mathit{oh}}
\def\Cohstack{\mathcal{C}\!\mathit{oh}}
\def\Spacesstack{\mathcal{S}\!\mathit{paces}}
\def\Quotfunctor{\mathrm{Quot}}
\def\Hilbfunctor{\mathrm{Hilb}}
\def\Curvesstack{\mathcal{C}\!\mathit{urves}}
\def\Polarizedstack{\mathcal{P}\!\mathit{olarized}}
\def\Complexesstack{\mathcal{C}\!\mathit{omplexes}}
% \Pic is the operator that assigns to X its picard group, usage \Pic(X)
% \Picardstack_{X/B} denotes the Picard stack of X over B
% \Picardfunctor_{X/B} denotes the Picard functor of X over B
\def\Pic{\mathop{\mathrm{Pic}}\nolimits}
\def\Picardstack{\mathcal{P}\!\mathit{ic}}
\def\Picardfunctor{\mathrm{Pic}}
\def\Deformationcategory{\mathcal{D}\!\mathit{ef}}
\usepackage{cleveref}
\crefname{originalreference}{original reference}{original references}
\Crefname{originalreference}{Original reference}{Original references}
\makeatletter
\newlabel{definition-relative-dualizing-proper-flat}{{0E5X}{source}{}{sourcepointer-1}{}}
\newlabel{definition-relative-dualizing-proper-flat@cref}{{[originalreference][0][]0E5X}{[1][0][]source}{}{}{}}
\makeatother
\crefname{originallemma}{original lemma}{original lemmas}
\Crefname{originallemma}{Original lemma}{Original lemmas}
\makeatletter
\newlabel{hypercovering-lemma-hypercovering-morphism-sites}{{0DAY}{source}{}{sourcepointer-2}{}}
\newlabel{hypercovering-lemma-hypercovering-morphism-sites@cref}{{[originallemma][0][]0DAY}{[1][0][]source}{}{}{}}
\makeatother
\crefname{originallemma}{original lemma}{original lemmas}
\Crefname{originallemma}{Original lemma}{Original lemmas}
\makeatletter
\newlabel{hypercovering-lemma-hypercovering-object}{{0DAV}{source}{}{sourcepointer-3}{}}
\newlabel{hypercovering-lemma-hypercovering-object@cref}{{[originallemma][0][]0DAV}{[1][0][]source}{}{}{}}
\makeatother
\crefname{originalremark}{original remark}{original remarks}
\Crefname{originalremark}{Original remark}{Original remarks}
\makeatletter
\newlabel{hypercovering-remark-take-unions-hypercovering-X}{{0DB2}{source}{}{sourcepointer-4}{}}
\newlabel{hypercovering-remark-take-unions-hypercovering-X@cref}{{[originalremark][0][]0DB2}{[1][0][]source}{}{}{}}
\makeatother
\crefname{originalsection}{original section}{original sections}
\Crefname{originalsection}{Original section}{Original sections}
\makeatletter
\newlabel{hypercovering-section-hypercovering-sites}{{094J}{source}{}{sourcepointer-5}{}}
\newlabel{hypercovering-section-hypercovering-sites@cref}{{[originalsection][0][]094J}{[1][0][]source}{}{}{}}
\makeatother
\crefname{originalreference}{original reference}{original references}
\Crefname{originalreference}{Original reference}{Original references}
\makeatletter
\newlabel{lemma-covering-enough}{{0E60}{source}{}{sourcepointer-6}{}}
\newlabel{lemma-covering-enough@cref}{{[originalreference][0][]0E60}{[1][0][]source}{}{}{}}
\makeatother
\crefname{originalreference}{original reference}{original references}
\Crefname{originalreference}{Original reference}{Original references}
\makeatletter
\newlabel{lemma-relative-dualizing-RHom}{{0E2G}{source}{}{sourcepointer-7}{}}
\newlabel{lemma-relative-dualizing-RHom@cref}{{[originalreference][0][]0E2G}{[1][0][]source}{}{}{}}
\makeatother
\crefname{originalreference}{original reference}{original references}
\Crefname{originalreference}{Original reference}{Original references}
\makeatletter
\newlabel{lemma-uniqueness-relative-dualizing}{{0E2D}{source}{}{sourcepointer-8}{}}
\newlabel{lemma-uniqueness-relative-dualizing@cref}{{[originalreference][0][]0E2D}{[1][0][]source}{}{}{}}
\makeatother
\crefname{originaldefinition}{original definition}{original definitions}
\Crefname{originaldefinition}{Original definition}{Original definitions}
\makeatletter
\newlabel{spaces-simplicial-definition-cartesian-derived-modules}{{0D9N}{source}{}{sourcepointer-9}{}}
\newlabel{spaces-simplicial-definition-cartesian-derived-modules@cref}{{[originaldefinition][0][]0D9N}{[1][0][]source}{}{}{}}
\makeatother
\crefname{originallemma}{original lemma}{original lemmas}
\Crefname{originallemma}{Original lemma}{Original lemmas}
\makeatletter
\newlabel{spaces-simplicial-lemma-fppf-glue-neg-ext-zero}{{0DLA}{source}{}{sourcepointer-10}{}}
\newlabel{spaces-simplicial-lemma-fppf-glue-neg-ext-zero@cref}{{[originallemma][0][]0DLA}{[1][0][]source}{}{}{}}
\makeatother
\crefname{originallemma}{original lemma}{original lemmas}
\Crefname{originallemma}{Original lemma}{Original lemmas}
\makeatletter
\newlabel{spaces-simplicial-lemma-fppf-neg-ext-zero-hom}{{0DL9}{source}{}{sourcepointer-11}{}}
\newlabel{spaces-simplicial-lemma-fppf-neg-ext-zero-hom@cref}{{[originallemma][0][]0DL9}{[1][0][]source}{}{}{}}
\makeatother
\title{Existence of the relative dualizing complex for a proper flat morphism}
\author{The Stacks Project Authors}
\date{stacks-project:spaces-duality.tex:local-snapshot-20261012; GNU FDL 1.2 or later}
\begin{document}
\maketitle
\noindent\textit{Editorial note.} All delivered mathematical text is the author's own, copied from the pinned Stacks Project chapter file spaces-duality.tex; no statement or proof step was rewritten, reordered or removed.\par\medskip
\noindent\textit{Editorial note.} The counted proof is the author's complete argument: Strong-flag prefilter hit re-read per the 2026-10-14 12:4x THIN-criterion patch (pb>=800 / pw>=100): the proof body was read in full at source; the main line is complete and the flagged phrase covers routine details, quoted verbatim in the editorial note below. Verbatim, proof line 2301: 'some details omitted -- compare with the end of the proof of Lemma ...' after the construction of the trace map; it covers the verification of the vanishing hypothesis, which is established locally and preserved by pullback.\par\medskip
\noindent\textit{Editorial note.} The label lemma-existence-relative-dualizing names the counted statement, and the proof environment immediately following it in the source is delivered verbatim as its proof.\par\medskip
\noindent\textit{Editorial note.} The counted statement is carried by the author's own lemma environment in the chapter source.\par\medskip
\noindent\textit{Editorial note.} Strong-flag prefilter hit re-read per the 2026-10-14 12:4x THIN-criterion patch: proof body (proof\_lines 2232-2314, pb=3537, pw=457) read in full; main line complete.\par\medskip
\noindent\textit{Editorial note.} Editorial note (verbatim): 'some details omitted -- compare with the end of the proof of Lemma ... for example to see why we have the required vanishing of negative exts.'\par\medskip
\noindent\textit{Editorial note.} Pre-open gate: python operations/agent-authoring/tools/precard\_gate.py --source spaces-duality.tex --label lemma-existence-relative-dualizing --start 2228 --end 2314 --local-id 28178 returned RESULT=OK. Cross-chapter same-name check: not carded in any other chapter.\par\medskip
\noindent\textit{Editorial note.} The statement and proof are copied from the pinned chapter file; the cited inputs are kept as external pointers with their canonical Stacks tags.\par\medskip
\noindent\textit{Allowed standard prerequisites (editorial).}\begin{itemize}
\item Etale hypercoverings with affine components and the induced hypercovering of the source. Hypercoverings, the lemmas producing a hypercovering whose components are affine and taking disjoint unions in each degree, and the lemma inducing a hypercovering of the source. Provides the simplicial algebraic spaces on which the relative dualizing complexes are constructed degree by degree.
\item Existence and uniqueness of relative dualizing complexes over affine bases. The discussion following the definition of a relative dualizing complex for a proper flat morphism of finite presentation over an affine base, together with the uniqueness lemma for pairs. Supplies the complexes on the affine components and the unique comparison isomorphisms needed for the gluing data.
\item Gluing objects and maps along hypercoverings with vanishing negative self extensions. Simplicial Spaces, the lemmas gluing cartesian derived modules and maps when negative self extensions vanish. Produces the global complex and the trace map from the simplicial gluing data.
\end{itemize}\medskip
\ifdefined\hypertarget\hypertarget{sourcepointer-1}{}\fi
\noindent\textit{Source reference.} Original reference 0E5X: sources/stacks-spaces-duality/pinned-original-source.tex, source lines 1978--1978. Prerequisite pointer invoked at proof line(s) 2251; it names a standard fact of the same chapter that the counted argument uses, and every step of the counted proof built on it is retained verbatim in the delivered body.\par\medskip
\ifdefined\hypertarget\hypertarget{sourcepointer-2}{}\fi
\noindent\textit{Source reference.} Original lemma 0DAY: sources/stacks-chow/tags.txt, source lines 16678--16678. Prerequisite pointer invoked at proof line(s) 2246; it names a standard fact of a sibling chapter that the counted argument uses, and every step of the counted proof built on it is retained verbatim in the delivered body.\par\medskip
\ifdefined\hypertarget\hypertarget{sourcepointer-3}{}\fi
\noindent\textit{Source reference.} Original lemma 0DAV: sources/stacks-chow/tags.txt, source lines 16675--16675. Prerequisite pointer invoked at proof line(s) 2239; it names a standard fact of a sibling chapter that the counted argument uses, and every step of the counted proof built on it is retained verbatim in the delivered body.\par\medskip
\ifdefined\hypertarget\hypertarget{sourcepointer-4}{}\fi
\noindent\textit{Source reference.} Original remark 0DB2: sources/stacks-chow/tags.txt, source lines 16682--16682. Prerequisite pointer invoked at proof line(s) 2240; it names a standard fact of a sibling chapter that the counted argument uses, and every step of the counted proof built on it is retained verbatim in the delivered body.\par\medskip
\ifdefined\hypertarget\hypertarget{sourcepointer-5}{}\fi
\noindent\textit{Source reference.} Original section 094J: sources/stacks-chow/tags.txt, source lines 11552--11552. Prerequisite pointer invoked at proof line(s) 2256; it names a standard fact of a sibling chapter that the counted argument uses, and every step of the counted proof built on it is retained verbatim in the delivered body.\par\medskip
\ifdefined\hypertarget\hypertarget{sourcepointer-6}{}\fi
\noindent\textit{Source reference.} Original reference 0E60: sources/stacks-spaces-duality/pinned-original-source.tex, source lines 2141--2141. Prerequisite pointer invoked at proof line(s) 2311; it names a standard fact of the same chapter that the counted argument uses, and every step of the counted proof built on it is retained verbatim in the delivered body.\par\medskip
\ifdefined\hypertarget\hypertarget{sourcepointer-7}{}\fi
\noindent\textit{Source reference.} Original reference 0E2G: sources/stacks-spaces-duality/pinned-original-source.tex, source lines 2054--2054. Prerequisite pointer invoked at proof line(s) 2267; it names a standard fact of the same chapter that the counted argument uses, and every step of the counted proof built on it is retained verbatim in the delivered body.\par\medskip
\ifdefined\hypertarget\hypertarget{sourcepointer-8}{}\fi
\noindent\textit{Source reference.} Original reference 0E2D: sources/stacks-spaces-duality/pinned-original-source.tex, source lines 2071--2071. Prerequisite pointer invoked at proof line(s) 2308; it names a standard fact of the same chapter that the counted argument uses, and every step of the counted proof built on it is retained verbatim in the delivered body.\par\medskip
\ifdefined\hypertarget\hypertarget{sourcepointer-9}{}\fi
\noindent\textit{Source reference.} Original definition 0D9N: sources/stacks-chow/tags.txt, source lines 16633--16633. Prerequisite pointer invoked at proof line(s) 2284; it names a standard fact of a sibling chapter that the counted argument uses, and every step of the counted proof built on it is retained verbatim in the delivered body.\par\medskip
\ifdefined\hypertarget\hypertarget{sourcepointer-10}{}\fi
\noindent\textit{Source reference.} Original lemma 0DLA: sources/stacks-chow/tags.txt, source lines 17040--17040. Prerequisite pointer invoked at proof line(s) 2286; it names a standard fact of a sibling chapter that the counted argument uses, and every step of the counted proof built on it is retained verbatim in the delivered body.\par\medskip
\ifdefined\hypertarget\hypertarget{sourcepointer-11}{}\fi
\noindent\textit{Source reference.} Original lemma 0DL9: sources/stacks-chow/tags.txt, source lines 17039--17039. Prerequisite pointer invoked at proof line(s) 2301; it names a standard fact of a sibling chapter that the counted argument uses, and every step of the counted proof built on it is retained verbatim in the delivered body.\par\medskip
\begin{lemma}
\label{lemma-existence-relative-dualizing}
Let $S$ be a scheme. Let $X \to Y$ be a proper, flat morphism of
algebraic spaces which is of finite presentation.
There exists a relative dualizing complex $(\omega_{X/Y}^\bullet, \tau)$.
\end{lemma}

\begin{proof}
Choose a \'etale hypercovering $b : V \to Y$ such that each
$V_n = \coprod_{i \in I_n} Y_{n, i}$ with $Y_{n, i}$ affine.
This is possible by Hypercoverings, Lemma
\ref{hypercovering-lemma-hypercovering-object} and
Remark \ref{hypercovering-remark-take-unions-hypercovering-X}
(to replace the hypercovering produced
in the lemma by the one having disjoint unions in each degree).
Denote $X_{n, i} = Y_{n, i} \times_Y X$ and $U_n = V_n \times_Y X$
so that we obtain an \'etale hypercovering
$a : U \to X$ (Hypercoverings, Lemma
\ref{hypercovering-lemma-hypercovering-morphism-sites})
with $U_n = \coprod X_{n, i}$.
For each $n, i$ there exists a relative dualizing complex
$(\omega_{n, i}^\bullet, \tau_{n, i})$ on $X_{n, i}/Y_{n, i}$.
See discussion following
Definition \ref{definition-relative-dualizing-proper-flat}.
For $\varphi : [m] \to [n]$ and $i \in I_n$ consider the morphisms
$g_{\varphi, i} : Y_{n, i} \to Y_{m, \alpha(\varphi)}$ and
$g'_{\varphi, i} : X_{n, i} \to X_{m, \alpha(\varphi)}$
which are part of the structure of the given hypercoverings
(Hypercoverings, Section \ref{hypercovering-section-hypercovering-sites}).
Then we have a unique isomorphisms
$$
\iota_{n, i, \varphi} :
(L(g'_{n, i})^*\omega_{n, i}^\bullet, Lg_{n, i}^*\tau_{n, i})
\longrightarrow
(\omega_{m, \alpha(\varphi)(i)}^\bullet, \tau_{m, \alpha(\varphi)(i)})
$$
of pairs, see discussion following
Definition \ref{definition-relative-dualizing-proper-flat}.
Observe that $\omega_{n, i}^\bullet$ has vanishing negative
self exts on $X_{n, i}$ by Lemma \ref{lemma-relative-dualizing-RHom}.
Denote $(\omega_n^\bullet, \tau_n)$ the pair on $U_n/V_n$ constructed using
the pairs $(\omega_{n, i}^\bullet, \tau_{n, i})$ for $i \in I_n$.
For $\varphi : [m] \to [n]$ and $i \in I_n$ consider the morphisms
$g_\varphi : V_n \to V_m$ and $g'_\varphi : U_n \to U_m$ which are part
of the structure of the simplicial algebraic spaces $V$ and $U$.
Then we have unique isomorphisms
$$
\iota_\varphi :
(L(g'_\varphi)^*\omega_n^\bullet, Lg_\varphi^*\tau_n)
\longrightarrow
(\omega_m^\bullet, \tau_m)
$$
of pairs constructed from the isomorphisms on the pieces.
The uniqueness guarantees that these isomorphisms satisfy
the transitivity condition as formulated in
Simplicial Spaces, Definition
\ref{spaces-simplicial-definition-cartesian-derived-modules}.
The assumptions of Simplicial Spaces, Lemma
\ref{spaces-simplicial-lemma-fppf-glue-neg-ext-zero}
are satisfied for $a : U \to X$, the complexes $\omega_n^\bullet$
and the isomorphisms $\iota_\varphi$\footnote{This lemma uses only
$\omega_0^\bullet$ and the two maps $\delta_1^1, \delta_0^1 : [1] \to [0]$.
The reader can skip the first few lines of the proof of the
referenced lemma
because here we actually are already given a simplicial system of
the derived category of modules.}.
Thus we obtain an object $\omega^\bullet$ of $D_\QCoh(\mathcal{O}_X)$
together with an isomorphism
$\iota_0 : \omega^\bullet|_{U_0} \to \omega_0^\bullet$
compatible with the two isomorphisms
$\iota_{\delta^1_0}$ and $\iota_{\delta^1_1}$.
Finally, we apply
Simplicial Spaces, Lemma
\ref{spaces-simplicial-lemma-fppf-neg-ext-zero-hom}
to find a unique morphism
$$
\tau : Rf_*\omega^\bullet \longrightarrow \mathcal{O}_Y
$$
whose restriction to $V_0$ agrees with $\tau_0$; some details omitted --
compare with the end of the  proof of
Lemma \ref{lemma-uniqueness-relative-dualizing}
for example to see why we have
the required vanishing of negative exts.
By Lemma \ref{lemma-covering-enough}
the pair $(\omega^\bullet, \tau)$ is a relative dualizing complex
and the proof is complete.
\end{proof}

\end{document}
